% !PRINTED 616
\weberpage{616}
%<<CONT
\weq{8}{\pi \equiv \beta \quad (\mathrm{mod}\ \mathfrak{m})}
\gesperrt{holds, and for which}
\weq{9}{\mathfrak{p} = \mathfrak{a}\pi}
\gesperrt{becomes a prime ideal.}

If one takes $\mathfrak{a} = 1$, then $\pi$ becomes an integer, and hence a prime existing in $\Omega$. For this particular case the theorem reads:

3. \gesperrt{In the field $\Omega$ there are infinitely many existing primes $\pi$ whose norm is a natural prime, and which are congruent, with respect to an arbitrary modulus $\mathfrak{m}$, to an arbitrary number $\beta$ coprime to $\mathfrak{m}$.}

If we take for the modulus $\mathfrak{m}$ an integer $\mu$ in $\Omega$, we obtain the following theorem:

4. \gesperrt{If $\mu$ is an integer in $\Omega$, then in the linear form}
\[
\mu\xi + \beta,
\]
\gesperrt{in which $\xi$ runs through the integers of $\Omega$ and $\beta$ is relatively prime to $\mu$, infinitely many primes are contained.}

Here one has a beautiful generalisation of the theorem on primes in arithmetic progressions (Bd.\ II, \S\,198)\footnote{For the case $\Delta = -4$ this proof has been given by \gesperrt{H. Weber}, Crelles Journ., Bd.\ 129 (Dirichlet-Band), for $\Delta = -3$ in a Strasbourg dissertation by \gesperrt{H. Bresslau} (Straßburg 1907).}.

\begin{center}
\textbf{\S\,168. Reduction of the class equation in the cyclotomic fields.}
\end{center}

We define a further group $A$ by the following determination:

Let $Q$, $m$ be two integral rational numbers whose prime factors we take up into the excludent $S$. Let $O$ be, as in \S\,98, the group of the numbers in $\Omega$, and $O'$ the group of the numbers of the order $[Q]$, and let $A$ consist of all integral and fractional numbers $\alpha$ in $O'$ for which the congruence condition
\weq{1}{N(\alpha) \equiv 1 \quad (\mathrm{mod}\ m)}
holds. First of all we have the class number
\weq{2}{h = (\overline{O}, \overline{E}A) = (\overline{O}, \overline{E}O')\,(\overline{E}O', \overline{E}A)}
%CONT>>
